Il settore agricolo è storicamente caratterizzato dalla necessità di ottimizzare le risorse, tra cui la gestione idrica rappresenta una delle criticità maggiori. In contesti di siccità prolungata, il razionamento dell'acqua durante i mesi estivi diventa una sfida determinante per la sopravvivenza delle colture. Parallelamente, l'impiego di fertilizzanti e pesticidi richiede un monitoraggio rigoroso, non solo per ragioni ambientali ma anche per garantire la conformità alle normative nazionali e comunitarie (UE), sempre più restrittive in termini di sostenibilità.

L'implementazione di reti di sensori, tecnologie GPS e sistemi di \textit{data collection} permette di ottimizzare radicalmente la gestione agricola. Queste soluzioni forniscono un monitoraggio continuo e in tempo reale dello stato fitosanitario delle colture, dei parametri \textit{edafici} (come \textbf{umidità} e \textbf{pH} del suolo) e dei pattern meteorologici, consentendo interventi mirati e tempestivi. 

Le metodologie di acquisizione dati spaziano dal monitoraggio on-site tramite reti di sensori e stazioni meteorologiche, fino all'impiego di immagini satellitari per le grandi estensioni agricole. L'elaborazione di questa mole di informazioni, attraverso algoritmi di \textbf{data analysis}, permette agli agricoltori di monitorare i tassi di crescita e implementare modelli previsionali per \textit{fitopatologie} e \textit{parassiti}, prevenendo perdite produttive ed economiche. Tali sistemi supportano inoltre il processo decisionale (decision making) riguardo alle tempistiche ottimali per semina, irrigazione e raccolta, estendendosi fino alla fase di post-raccolta per il monitoraggio dei parametri di stoccaggio, come nel caso dei cereali. \cite{iot_web_service_agricolture}\newline

\noindent L'esempio più emblematico di innovazione nel settore è l' \textbf{agricoltura di precisione}, dove l'integrazione di ecosistemi IoT permette di monitorare costantemente i parametri edafici e atmosferici. Tale approccio consente di calibrare con estrema accuratezza l'apporto di input agronomici quali acqua, fertilizzanti e fitofarmaci, massimizzando l'efficienza delle risorse e riducendo l'impatto ambientale.

Tra le applicazioni più rilevanti dell'IoT in ambito agricolo si distinguono:

\begin{itemize}
	\item \textbf{Irrigazione di precisione}: consente la distribuzione idrica mirata, attivandola solo in risposta al reale stress idrico della pianta. Questo previene gli sprechi e ottimizza il mantenimento del punto di appassimento critico del terreno.
	\item \textbf{Monitoraggio post-raccolto}: i cereali costituiscono la base della sicurezza alimentare globale; monitorarne lo stato dopo la raccolta è fondamentale per prevenire la proliferazione di muffe o micotossine, garantendo l'integrità del prodotto destinato al consumo umano.
	\item \textbf{Smart Greenhouse}: l'impiego di sensori e attuatori IoT permette di automatizzare la gestione dei microclimi, rendendo possibile la coltivazione di specie esigenti anche in periodi dell'anno climaticamente sfavorevoli, attraverso il controllo rigoroso di temperatura, umidità e radiazione solare.
	\item \textbf{Monitoraggio della salute delle colture}: l'impiego di droni (UAV) equipaggiati con sensori multispettrali permette di sorvegliare vaste estensioni di terreno, identificando tempestivamente anomalie vegetative, carenze nutrizionali o attacchi parassitari non visibili a occhio nudo.
\end{itemize}

\noindent Un ulteriore ambito di rilievo, strettamente connesso alla digitalizzazione rurale, è il \textbf{monitoraggio del bestiame}. Attraverso l'adozione di sensori \textit{wearable}, è infatti possibile tracciare in tempo reale lo stato di salute e i pattern comportamentali degli animali, migliorando sensibilmente il benessere zootecnico e la resa produttiva \cite{iot_agricultural_7realexamples}\cite{digi_iot_agricolture}.